\noindent\textbf{Overview.} We argue by contradiction and slice by Haar-random subspaces $V$ (Lemma~\ref{ang:slice-selection}). If $q>k-1$, planar slices ($h=k-1$) already contradict the planar obstruction (Lemma~\ref{ang:planar-measures}). If $q=k-1$, planar slices would leave direction exponent $0$. So we slice to dimension three ($h=k-2$), where the point exponent is $t-h>3/2$ and the direction exponent is $1$. Three tube covers in transverse directions confine a set to about $\delta^{-3/2}$ balls (Lemma~\ref{ang:loomis-whitney}), and this contradicts the point ball bound. The real work is finding, inside one slice, three \emph{quantitatively transverse} directions that share a large incident point set. Here the plate hypothesis~\eqref{ang:endpoint-plate} enters, through the weak radial estimate.

\begin{proof}
\textbf{Step 0: Reductions.} It suffices to rule out a sequence $\delta\to0$ of data obeying the hypotheses with losses $\delta^{-o(1)}$, tube counts $\delta^{-q-o(1)}$ and $(\mu\otimes\nu)(E)\ge\delta^{o(1)}$. Indeed, if no $\epsilon,c$ worked, counterexamples with $\epsilon=c=1/j$ at scales $\delta_j<1/j$ would form such a sequence. Cover $S^k$ by finitely many small caps and pigeonhole one carrying incidence $\ge\delta^{o(1)}$. After a rotation, which preserves every hypothesis, this cap is $w(B)$, where $B\subset\R^k$ is a small ball and $w(p)=(1,p)/\sqrt{1+|p|^2}$. The map $w$ is bi-Lipschitz on $B$. Restrict $\nu$ and $E$ to this cap and normalize $\nu$. Its mass is $\ge\delta^{o(1)}$ (it bounds the incidence), so all losses stay subpower.

A selected slice is always normalized in the same way. Lemma~\ref{ang:slice-selection} gives $\mathbb E_V\int(\widetilde\mu_{V,z}\otimes\widetilde\nu_V)(\widetilde E_{V,z})\,dz\ge\delta^{o(1)}$ over a bounded $z$-domain. Hence some $(V,z)$ has incidence $I\ge\delta^{o(1)}$. Both masses are $\le\delta^{-o(1)}$ and $I$ is at most their product, so both masses are also $\ge\delta^{o(1)}$. Dividing by them gives probabilities with subpower losses and incidence $\ge\delta^{o(1)}$.

\textbf{Step 1: The case $q>k-1$.} Apply Lemma~\ref{ang:slice-selection} with $h=k-1$; its requirements $t,q>h$ and $h\le n-2=k-1$ hold. Then $V$ is a plane, and
\[
 b=q-h\in(0,\tfrac12],\qquad f=t-h\in(\tfrac12,1),\qquad b<f
\]
because $q\le k-\tfrac12<t$. Take a normalized slice from Step~0 and use the original points and directions as labels, with coordinate maps $p_{V,z}$ and $v_V$. Item~3 gives the $f$- and $b$-ball bounds of the pushforwards, and for each label $w$ a cover by $\delta^{-b-o(1)}$ strips parallel to $v_V(w)$. Along our sequence this is exactly the data excluded by Lemma~\ref{ang:planar-measures}.

\textbf{Step 2: Setup for $q=k-1$.} Put $h=k-2$, so $\dim V=3$, and fix
\[
 g=t-h-\tfrac32>0,\qquad 0<a<g/4,\qquad h<\ell<h+\tfrac1k=\tfrac{(k-1)^2}{k}.
\]
Let $\lambda$ be the normalized law of $p=w^{-1}(\omega)$, $\omega\sim\nu$, and let $\nu=w_\#\lambda$ be what we feed to slicing. Then $\lambda$ has $(k-1)$-ball bounds with subpower loss. It also satisfies the affine plate bound $\lambda(H^{(r)})\lesssim\delta^{-o(1)}r^\zeta$, for the following reason. An affine hyperplane $\{a_0+a\cdot p=0\}$ is mapped by $w$ into the linear hyperplane $\{a_0y_0+a\cdot y=0\}$, and $w$ is Lipschitz, so~\eqref{ang:endpoint-plate} applies. Write $\lambda_V=\delta^{-h}\one_{\{\dist(w(p),V)\le\delta\}}\lambda$. By Fubini,
\[
 \mathbb E_V\lambda_V^{\otimes j}(S)=\int_S K_j\,d\lambda^{\otimes j},\qquad
 K_j=\delta^{-jh}\Pr_V\{\dist(w(p_i),V)\le\delta\ \forall i\}.
\]

\emph{Wedge geometry.} Put $u_i=p_i-p_0$. Then $(1,p_0)\wedge(1,p_1)\wedge\cdots=(1,p_0)\wedge(0,u_1)\wedge\cdots$. The part containing $e_0$ is orthogonal to the rest, which gives
\[
|(1,p_0)\wedge(0,u_1)\wedge\cdots\wedge(0,u_m)|^2=|u_1\wedge\cdots\wedge u_m|^2+|p_0\wedge u_1\wedge\cdots\wedge u_m|^2 .
\]
Since $p_0$ is bounded and normalization costs constants, this yields
\begin{equation}\label{wedge}
 |w(p_0)\wedge w(p_1)|\asymp|u_1|,\qquad |w(p_0)\wedge w(p_1)\wedge w(p_2)|\asymp|u_1||u_2|\,\alpha,
\end{equation}
where $\alpha=|\tfrac{u_1}{|u_1|}\wedge\tfrac{u_2}{|u_2|}|$. For unit columns the top singular value is $\asymp1$, and the products of the leading singular values are the wedges in~\eqref{wedge}. So Lemma~\ref{ang:haar} (with $n=k+1$, $j\le3=n-h$) gives
\begin{equation}\label{K2}
K_2(p_0,p_1)\le C(|u_1|+\delta)^{-h}.
\end{equation}

\textbf{Step 3: Exceptional pairs.} Let $e>0$ be small and fixed (chosen in Step~4). Lemma~\ref{ang:weak-radial} applies to $\lambda$ at small scales, since its input losses are subpower. It gives $c_e>0$ and a pair set $G$ with $(\lambda\otimes\lambda)(G^c)\lesssim\delta^{c_e}$ and the cap bound with exponent $\ell$ and loss $\delta^{-e}$. For $b_0>0$ (Step~4) set
\[
\mathcal E_{\rm pair}=G^c\cup\{|p_0-p_1|<\delta^{b_0}\}.
\]
Call an anchored pair \emph{retained} if it lies outside $\mathcal E_{\rm pair}$.

\emph{Claim: $\mathbb E_V(\lambda_V\otimes\lambda_V)(\mathcal E_{\rm pair})$ is power-small.} By~\eqref{K2} and the $(k-1)$-ball bound, summing over dyadic annuli $|u_1|\sim r$ gives
\[
\int_{|p_0-p_1|\le s}K_2\,d\lambda^{\otimes2}\le\delta^{-o(1)}\sum_{\delta\le r\le s}r^{k-1-h}\le\delta^{-o(1)}s .
\]
Here $k-1-h=1$, and the term $|u_1|<\delta$ is included at $r=\delta$. Fix $\theta\in(0,1]$ with $\theta h<c_e$. On $G^c\cap\{|u_1|>\delta^\theta\}$ we have $K_2\le C\delta^{-\theta h}$. Hence
\[
\mathbb E_V(\lambda_V\otimes\lambda_V)(\mathcal E_{\rm pair})\le C\delta^{c_e-\theta h}+\delta^{\theta-o(1)}+\delta^{b_0-o(1)}.
\]

\textbf{Step 4: Nearly dependent triples.} Let $\mathcal T$ be the set of triples $(p_0,p_1,p_2)$ such that both pairs $(p_0,p_i)$ are retained and $|w(p_0)\wedge w(p_1)\wedge w(p_2)|<2\delta^a$. On $\mathcal T$, \eqref{wedge} and $|u_i|\ge\delta^{b_0}$ give $\alpha\le R:=C\delta^{a-2b_0}$. They also give singular values $d_1\asymp1$, $d_2\gtrsim\delta^{b_0}$ and $d_3\gtrsim\delta^{2b_0}\alpha$, so Lemma~\ref{ang:haar} gives
\[
K_3\le C\delta^{-3hb_0}(\alpha+\delta)^{-h}.
\]
Fix $p_0,p_1$. The condition $\alpha\le r$ puts the displacement direction of $p_2$ from $p_0$ in two antipodal caps of radius $\asymp r$. So the weak radial estimate gives $\lambda\{p_2:(p_0,p_2)\in G,\ \alpha\le r\}\le C\delta^{-e}r^\ell$. Summing dyadic shells $\alpha\sim r$, $\delta\le r\le R$, together with the shell $\alpha\le\delta$, and using $\ell>h$, gives
\begin{equation}\label{badT}
\mathbb E_V\lambda_V^{\otimes3}(\mathcal T)\le C\delta^{-3hb_0-e}R^{\ell-h}\le C\delta^{(a-2b_0)(\ell-h)-3hb_0-e}.
\end{equation}
\emph{Order of constants:} $a$ and $\ell$ are fixed above. Next choose $b_0$ small so that $a>2b_0$ and $(a-2b_0)(\ell-h)>3hb_0$. Then choose $e$ small so that the exponent in~\eqref{badT} is positive. This fixes $c_e$, and then $\theta$. All of these are independent of $\delta$.

\textbf{Step 5: A good triple in one slice.} Apply Lemma~\ref{ang:slice-selection} with $h=k-2$; again $t,q>h$ holds. Write $\widetilde\lambda_V=(w^{-1})_\#\widetilde\nu_V\le\lambda_V$; its mass is at most $\delta^{-o(1)}$ uniformly in $V$. For each $V$,
\[
\int\|\mu_{V,z}\|\,dz=\int\!\!\int\delta^{-h}\one_{\{|P_{V^\perp}x-z|\le\delta\}}\,dz\,d\mu(x)\le C .
\]
So the measure $dm=d\widetilde\mu_{V,z}(x)\,dz\,dV$ has total mass $\le C$. Let
\[
F(V,z,x)=\widetilde\lambda_V\{p:(x,w(p))\in\widetilde E_{V,z}\}.
\]
Then $\int F\,dm\ge\delta^{o(1)}$, and H\"older gives $\int F^3\,dm\ge\delta^{o(1)}$. Now $\int F^3\,dm$ is the $m\otimes\widetilde\lambda_V^{\otimes3}$-mass of tuples $(V,z,x,p_0,p_1,p_2)$ with $(x,w(p_i))\in\widetilde E_{V,z}$ for all $i$. Drop the incidence constraint and use $\widetilde\lambda_V\le\lambda_V$ and the uniform bound $\int\|\widetilde\mu_{V,z}\|dz\le C$; no independence of the point and direction slices is needed. Then the tuples with $(p_0,p_1,p_2)\in\mathcal T$ have mass $\le C\,\mathbb E_V\lambda_V^{\otimes3}(\mathcal T)$. Those with $(p_0,p_i)\in\mathcal E_{\rm pair}$ have mass
\[
\le C\sup_V\|\widetilde\lambda_V\|\;\mathbb E_V(\lambda_V\otimes\lambda_V)(\mathcal E_{\rm pair}).
\]
Both are power-small by~\eqref{badT} and Step~3. So mass $\ge\delta^{o(1)}$ remains on triples with wedge $\ge2\delta^a$. The parameters $(V,z,p_0,p_1,p_2)$ carry total mass $\le\delta^{-o(1)}$. So we can fix one slice $(V,z)$ and one such triple whose common incident set
\[
A_0=\{x:(x,w(p_i))\in\widetilde E_{V,z},\ i=0,1,2\}
\]
has $\widetilde\mu_{V,z}(A_0)\ge\delta^{o(1)}$.

\textbf{Step 6: Contradiction.} Let $A=p_{V,z}(A_0)\subset V\cong\R^3$. By item~3 the pushforward of $\widetilde\mu_{V,z}$ has $(t-h)$-ball bounds with subpower loss, and it gives $A$ mass $\ge\delta^{o(1)}$. Hence
\[
|A|_\delta\ge\delta^{-(t-h)+o(1)}.
\]
For each $i$, item~3 also covers $A$ by $\delta^{-(q-h)-o(1)}=\delta^{-1-o(1)}$ tubes of width $C'\delta$ parallel to $v_i=v_V(w(p_i))$. Since $\dist(w(p_i),V)\le\delta$, we have $|v_i-w(p_i)|=O(\delta)$, so $|\det(v_0,v_1,v_2)|\ge2\delta^a-O(\delta)\ge\delta^a$. Lemma~\ref{ang:loomis-whitney} with $m=3$ and $\Delta=\delta^a$ gives
\[
|A|_\delta\le C\delta^{-3a}\big(\delta^{-1-o(1)}\big)^{3/2}=\delta^{-3/2-3a-o(1)} .
\]
Since $3a<g=t-h-\tfrac32$, the two bounds on $|A|_\delta$ are incompatible for small $\delta$.

When $k=2$ we have $h=0$. Then $V=\R^3$, all slicing weights and Haar factors equal $1$, the $z$-integral is over a point, and the argument is unchanged.
\end{proof}
